% Part: first-order-logic
% Chapter: beyond
% Section: introduction

\documentclass[../../../include/open-logic-section]{subfiles}

\begin{document}

\olfileid{fol}{byd}{int}

\olsection{لمحة عامة}

لا ينحصر المنطق الجدير بالنظر في منطق الرتبة الأولى؛ فله وجوه
كثيرة من التوسعة والتغيير. والنسق المنطقي، في الغالب، تعيين صوري
للغة، ومعه عادةً نسق للاستنباط ودلالة مقصودة، وإن لم يجتمعا في
كل نسق. غير أن تسمية هذه الأنساق «منطقًا» بالاصطلاح التقني
تثير سؤالًا: ما صلتها بالمنطق في الفهم الحدسي أو الفلسفي؟
فالأنساق الآتية كلها موضوعة لتمثيل ضرب من الاستدلال؛ فما الذي
يجعل ذلك الاستدلال منطقيًا؟

ليس لهذا السؤال جواب يسير. فكلمة «منطق» تختلف معانيها ومواضع
استعمالها، وقد تناول الفلاسفة مفهومها من جهات شتى، كما تناولوا
مفهوم «الصدق». فقد يُجعل غرض الاستدلال المنطقي معرفة ما يصدق
بالضرورة، أو ما يصدق قبليًا، أو ما لا يتوقف صدقه على تأويل
الألفاظ غير المنطقية، أو ما يصدق بصورته، أو بالعرف اللغوي.
وكل واحد من هذه المعاني يحتاج إلى بيان كثير. بل لو قصرنا النظر
على المنطق المستعمل في الرياضيات، لم يرتفع الخلاف في حدوده.
ففي \textit{برنسيبيا ماثماتيكا} حاول راسل ووايتهد تأسيس الرياضيات
على المنطق، على نهج {\em المنطقيانية} الذي ابتدأه فريغه.
وكان نسقهما منطقًا للأنماط العليا، شبيهًا بما سيأتي وصفه.
ثم اضطرا إلى إدخال بديهيات لا تُعد منطقية خالصة في أكثر المعايير،
ولا سيما بديهية اللانهاية وبديهية القابلية للرد. ومع ذلك قد يرى
الناظر أن بعض الاستدلال من الرتب العليا ينبغي عدّه منطقيًا.
وأما كواين، الذي لا يجيز في أنطولوجيته «القضايا» موضوعات للخطاب،
فيرى خلاف ذلك: منطق الرتبة الثانية والرتب العليا عنده نظرية
مجموعات متسترة بلباس المنطق؛ أي إن التكميم على المحمولات ليس
عنده من المنطق الخالص.

نترك الآن الفصل في هذه المسائل الفلسفية، وننظر في الأنساق
التالية على أنها تمثيلات صورية مثالية لضروب من الاستدلال،
منطقيّة كانت أو غير منطقيّة.

\end{document}
